%% ch24_universal_portability.tex
%% Universal Portability: Formal Argument and Evidence
%% Structured claim with six-domain evidence base.

\subsection{Universal Portability: Formal Argument}
\label{sec:universal_portability}

This section formalises the claim that DC3S is \emph{universally portable}:
any cyber-physical system (CPS) that satisfies four mild structural conditions
can use the DC3S safety layer without domain-specific modification to the core
pipeline.

% ------------------------------------------------------------------
\subsubsection*{Structural Conditions for Portability}

\begin{definition}[DC3S-Compatible CPS]
\label{def:compatible_cps}
A cyber-physical system $\mathcal{P} = (\mathcal{S}, \mathcal{A}, \phi, \tau)$
is \emph{DC3S-compatible} if:
\begin{enumerate}[label=(C\arabic*)]
  \item \textbf{Observable state.} There exists a telemetry function
        $\psi: \mathcal{S}_\text{true} \to \mathbb{R}^m$ such that
        $s_t^\text{obs} = \psi(s_t^\text{true}) + \varepsilon_t$
        with $\varepsilon_t$ bounded in expectation.
  \item \textbf{Numeric action space.} The action space $\mathcal{A} \subseteq \mathbb{R}^d$
        is a non-empty bounded interval per dimension.
  \item \textbf{Expressible safety predicate.} The safety predicate
        $\phi: \mathcal{S} \to \{0,1\}$ can be written as a finite conjunction of
        box constraints: $\phi(s) = \mathbf{1}[\forall i: \ell_i \leq s_i \leq u_i]$.
  \item \textbf{Calibration data.} At least $n_\text{cal} \geq 20$ historical
        (state, action, outcome) tuples are available for conformal calibration.
\end{enumerate}
\end{definition}

Conditions (C1)--(C3) are satisfied by the vast majority of engineered CPS:
they exclude only systems with purely discrete action spaces or unobservable
states.  Condition (C4) is a minimal data requirement met by any deployed system
after a brief warm-up period.

% ------------------------------------------------------------------
\subsubsection*{Portability Claim}

\begin{claim}[Universal DC3S Portability]
\label{claim:portability}
For any DC3S-compatible CPS $\mathcal{P}$, the five-stage DC3S pipeline
(Detect $\to$ Calibrate $\to$ Constrain $\to$ Shield $\to$ Certify) runs
without modification, and the following guarantees hold:
\begin{enumerate}[label=(\roman*)]
  \item \emph{Marginal coverage}: $\mathbb{E}[\text{TSVR}] \leq \alpha(1 - \bar{w})$
        (Theorem~1).
  \item \emph{Group-conditional coverage}: PICP $\geq 1-\alpha$ per reliability
        group (Theorem~9).
  \item \emph{High-probability bound}: $P(\text{TSVR} \geq \alpha(1-\bar{w}) + \varepsilon)
        \leq e^{-2T\varepsilon^2}$ (Theorem~10).
  \item \emph{Repair feasibility}: $a_\text{safe} \in [\ell_t, u_t]$ always
        (Proposition~1).
  \item \emph{Runtime overhead}: p99 latency $< 0.35$\,ms per step
        (empirical, Table~\ref{tbl:latency_benchmark}).
\end{enumerate}
\end{claim}

\textbf{Proof sketch.}
Guarantees (i)--(iii) follow from Theorems~1, 9, 10, which are stated for
generic $w_t \in [0,1]$ and $\phi$ satisfying (C3) — no domain physics enters
the proof.  Guarantee (iv) holds by algebraic construction of the clamping
repair (Proposition~1 + condition (C2)).  Guarantee (v) is bounded by
$O(W \log W + d)$ per step (Table~\ref{tbl:complexity_analysis}), which for
typical $W=20$, $d \leq 6$ is sub-millisecond on any modern CPU. $\square$

% ------------------------------------------------------------------
\subsubsection*{Six-Domain Evidence Base}

Claim~\ref{claim:portability} is validated across six CPS domains spanning
three orders of magnitude in action space dimensionality, three physical
modalities (energy, mobility, biology/chemistry), and two data regimes (synthetic
ORIUS-Bench and real UCI/PhysioNet data):

\begin{center}
\begin{tabular}{llll}
\toprule
\textbf{Domain} & \textbf{Modality} & \textbf{Data} & \textbf{TSVR Reduction} \\
\midrule
Battery  & Energy storage & DE/US grid (locked) & Reference \\
Vehicle  & Ground mobility & Synthetic kinematic & 96.9\% \\
Healthcare & Biological CPS & PhysioNet BIDMC & 72.1\% \\
Industrial & Thermal process & UCI CCPP & 100\% \\
Aerospace & Aerial vehicle & Synthetic flight env. & 100\% \\
Navigation & Robotic 2-D & Synthetic arena & 100\% \\
\bottomrule
\end{tabular}
\end{center}

All five proof domains exceed the 25\% evidence gate (Table~\ref{tbl:multi_domain_evidence_gate}).

% ------------------------------------------------------------------
\subsubsection*{Scope Limitation and Future Work}

Claim~\ref{claim:portability} does \emph{not} assert:
\begin{itemize}
  \item \textbf{Optimality} of the repair.  DC3S minimises conservatism by
        inflating only proportional to $w_t$, but does not solve a full
        constrained optimisation at each step.  A model-predictive repair
        layer could reduce intervention burden further at higher computational cost.
  \item \textbf{Convergence} of the controlled system.  DC3S is a one-step
        safety filter; closed-loop stability of the underlying controller plus
        DC3S is a separate systems-theoretic question outside the scope of this thesis.
  \item \textbf{Hardware-validated deployment} in any non-battery domain.
        The vehicle, healthcare, industrial, aerospace, and navigation results
        use ORIUS-Bench synthetic dynamics.  Field pilots are identified as
        future work in \S\ref{sec:future_work}.
\end{itemize}

These limitations are honest and deliberate: the universal portability claim
is a \emph{framework contribution}, not an operational deployment claim.  The
thesis establishes the theoretical and empirical conditions under which DC3S
transfers; translating that to production in each domain requires domain-specific
engineering work beyond the scope of a single dissertation.
